% !TeX root = ../../Book5.tex
% 纲要标题：### 9.2 表示与模
% 纲要位置：工作记录/计划纲要.md，第 4017 行。

\section{表示与模}
\label{b5:sec:0902}


% 纲要范围：
%
% * 向量空间、线性箭及交换关系
% * 箭图表示与路径代数模的等价
% * 左模约定与路径复合方向
% * 以菱形箭图核验换基不变量及关系对同构参数的限制
%

路径代数把箭的复合变成乘法，箭图表示则把它变成线性映射的复合。本节不仅比较两者的对象，还比较保持作用的映射；这使直和、核、商和扩张都能在两种描述之间转换。

\subsection{向量空间、线性箭及交换关系}
\label{b5:subsec:0902:01}

\begin{definition}\label{b5:def:quiver-representation}
设 \(k\) 为域，\(Q\) 为有限箭图。给每个顶点 \(i\) 一个有限维 \(k\) 向量空间 \(V_i\)，给每支箭 \(a:i\to j\) 一个线性映射 \(V_a:V_i\to V_j\)，所得数据 \(V=(V_i,V_a)\) 称为 \(Q\) 的有限维 \(k\)\term[jiantu-biaoshi]{表示}[representation of a quiver]。对路径 \(p=a_r\cdots a_1\)，记 \(V_p=V_{a_r}\cdots V_{a_1}\)，并令 \(V_{e_i}=1_{V_i}\)。
对两个表示 \(V,W\)，一族线性映射 \(f_i:V_i\to W_i\) 若满足
\[
W_a f_i=f_j V_a\qquad(a:i\to j),
\]
就称为表示同态 \(f:V\to W\)。若所有 \(f_i\) 都可逆，则称 \(f\) 为表示同构。
\end{definition}
对于有关系箭图 \((Q,I)\)，再要求每个关系 \(\sum_p c_p p\in I\) 所给的映射 \(\sum_p c_pV_p\) 为零。只需检查一组双边理想生成关系：左右乘路径就是在两侧复合映射，加法则是相加。记这些表示及其同态组成的范畴为 \(\operatorname{rep}_k(Q,I)\)；\(I=0\) 时写为 \(\operatorname{rep}_k Q\)。
\symidx{\(\operatorname{rep}_k(Q,I)\)：有关系箭图的有限维表示范畴}
\begin{definition}\label{b5:def:quiver-subrepresentation}
设 \(k\) 为域，\(Q\) 为有限箭图，\(V=(V_i,V_a)\)、\(W=(W_i,W_a)\) 为 \(Q\) 的有限维 \(k\) 表示。若一族子空间 \(U_i\subseteq V_i\)、\(i\in Q_0\)，满足 \(V_a(U_i)\subseteq U_j\) 对每支箭 \(a:i\to j\) 成立，则这些子空间及箭映射的限制组成 \(V\) 的子表示 \(U\)。商表示在各顶点取 \(V_i/U_i\)，箭映射由 \(V_a\) 诱导。直和 \(V\oplus W\) 则在各顶点及各支箭分别取直和。向量 \(\mathbf d(V)=(\dim_kV_i)_{i\in Q_0}\) 称为 \(V\) 的\term[weishu-xiangliang]{维数向量}[dimension vector]。
\end{definition}
\symidx{\(\mathbf d(V)\)：箭图表示的维数向量}
% 跨卷引用目标：b4:def:abelian-category（第四卷定义4.2.1）。
表示同态的核与余核也逐顶点计算。例如 \(W_a f_i=f_jV_a\) 使 \(V_a\) 将 \(\ker f_i\) 映入 \(\ker f_j\)，而同一个等式保证商上的箭映射良定义。因此一列表示正合，当且仅当它在每个顶点的向量空间列都正合。这也给出第四卷定义4.2.1所说的交换范畴结构：同态逐顶点相加，直和为有限双积，核与余核均存在；逐顶点的第一同构定理又使余像到像的自然映射为同构。因此 \(\operatorname{rep}_k(Q,I)\) 是交换范畴。

\subsection{箭图表示与路径代数模的等价}
\label{b5:subsec:0902:02}

\begin{theorem}\label{b5:thm:quiver-module-equivalence}
设 \(k\) 为域，\(Q\) 为非空有限箭图，\(I\) 为 \(kQ\) 的双边理想，\(e_i\) 为顶点 \(i\) 的长度零路径。有限维左 \(kQ/I\) 模范畴与满足 \(I\) 中全部关系的有限维 \(k\) 箭图表示范畴 \(\operatorname{rep}_k(Q,I)\) 等价。对应在对象上为
\[
M\longmapsto(e_iM,\ a|_{e_iM}),\qquad
(V_i,V_a)\longmapsto\bigoplus_{i\in Q_0}V_i.
\]
第二个对象上，顶点幂等元作坐标投影，路径按相应线性映射的复合作用；两种对应保持直和与正合列。
\end{theorem}
\begin{proof}
若 \(M\) 是左模，单位分解为 \(\sum_i e_i\)，故 \(M=\bigoplus_i e_iM\)。由 \(a=e_jae_i\)，左乘 \(a\) 将 \(e_iM\) 映入 \(e_jM\)，而 \(I\) 在 \(M\) 上作用为零，得到满足关系的表示。模同态把每个 \(e_iM\) 送入目标模的相应分量，并与每条箭对应的线性映射相容，因而给出表示同态。
反过来，给定表示，在直和上令一条从 \(i\) 到 \(j\) 的路径只作用于第 \(i\) 个坐标，并用 \(V_p\) 将其送到第 \(j\) 个坐标。如果两条路径不接合，两个坐标算子的复合为零；如果接合，复合正是拼接路径的算子。因此这些算子延拓为 \(kQ\) 的左作用。关系条件使作用通过 \(kQ/I\)。一族表示同态的直和保持顶点分量，并与两端的箭映射相容，所以也与全部路径作用相容，成为模同态。
从模出发再取直和，由 \((m_i)_i\mapsto\sum_i m_i\) 得到与作用相容的自然同构；从表示出发则在各坐标恢复原空间和原箭映射。以上构造对同态也互逆，故为范畴等价。直和逐坐标保持，正合性由各顶点的核与余核计算保持。
\end{proof}
例如单圈表示就是一个向量空间及其一个线性算子，因而等同于 \(k[t]\) 模；两个圈的表示是两个任意线性算子，不要求它们交换。若另加关系 \(ab-ba=0\)，才变为两个可交换算子。箭图与关系分别记录“有哪些映射”和“映射应满足什么恒等式”。

\ProseParagraphBreak

路径在某一个表示上作用为零，与它在路径代数中为零是两件事。在 \(Q=(1\xrightarrow{a}2\xrightarrow{b}3)\) 中，\(ba\) 是 \(kQ\) 的一个非零基元素。取表示
\[
V=\bigl(k\xrightarrow{1}k\xrightarrow{0}k\bigr),
\]
则 \(V_bV_a=0\)，说明 \(ba\) 落入这个表示的作用同态
\(kQ\to\operatorname{End}_k(V_1\oplus V_2\oplus V_3)\) 的核，并未改变 \(kQ\) 本身。如果改研究 \(B=kQ/(ba)\)，则每个 \(B\) 模都必须满足这个零复合条件；例如 \(k\xrightarrow{1}k\xrightarrow{1}k\) 是 \(kQ\) 的表示，却不是 \(B\) 的表示。取商是在所有表示中共同规定关系，单个表示的核则可以额外零化更多路径。

\subsection{左模约定与路径复合方向}
\label{b5:subsec:0902:03}

本章始终用左模。若 \(a:i\to j\)，左乘 \(a=e_jae_i\) 从 \(e_iM\) 到 \(e_jM\)，与箭向相同。若改用右模，则右乘 \(a\) 从 \(Me_j\) 到 \(Me_i\)，箭向反转。因此把一个右模版本直接抄成左模，往往会同时弄反路径乘积、投射模和扩张箭的方向。
\ProseParagraphBreak
一个能随时检查约定的例子是上三角矩阵代数 \(T_2(k)\)。令 \(e_1=E_{11}\)、\(e_2=E_{22}\)，则 \(E_{12}=e_1E_{12}e_2\)，对应的箭为 \(2\to1\)。把箭送到 \(E_{12}\)，顶点送到两个对角矩阵，便给出此箭图的路径代数与 \(T_2(k)\) 的同构：两侧均以这三个元素为基，全部乘法关系相同。
\begin{proposition}\label{b5:prop:quiver-change-of-basis}
设 \(k\) 为域，\(Q\) 为有限箭图，\(\mathbf d=(d_i)_{i\in Q_0}\) 为非负整数向量。固定每个顶点空间为 \(k^{d_i}\)，则所有 \(k\) 表示由矩阵组
\[
\operatorname{Rep}_k(Q,\mathbf d)
=\prod_{a:i\to j}\operatorname{Mat}_{d_j\times d_i}(k)
\]
给出。换基群 \(\prod_i\operatorname{GL}_{d_i}(k)\) 作用为
\[
(g_i)_i\cdot(X_a)_a=(g_jX_ag_i^{-1})_{a:i\to j}.
\]
两个矩阵组在同一轨道中，当且仅当相应表示同构。加关系时，取满足关系矩阵方程的稳定子集即可。
\end{proposition}
\symidx{\(\operatorname{Rep}_k(Q,\mathbf d)\)：固定维数向量的箭图表示矩阵空间}
\begin{proof}
每支箭的线性映射由任意相应尺寸的矩阵给定，故得到所列直积。表示同构的方程为 \(Y_ag_i=g_jX_a\)，每个 \(g_i\) 可逆时恰好等价于 \(Y_a=g_jX_ag_i^{-1}\)。关系的矩阵值从起点到终点作同样的左右换基，因而其为零的条件在作用下保持。
\end{proof}
维数向量只规定矩阵的尺寸，并不确定表示。例如一个圈上的一维表示由标量 \(\lambda\in k\) 给出，一维换基不改变这个标量，所以不同的 \(\lambda\) 已经给出不同构的表示。即使没有有向环，也可能留下不能由换基消去的参数。

\begin{proposition}\label{b5:prop:diamond-scalar-representations}
设 \(k\) 为域，\(Q\) 有顶点 \(1,2,3,4\) 及箭
\(a:1\to2\)、\(b:1\to3\)、\(c:2\to4\)、\(d:3\to4\)。在每个顶点放置 \(k\)，在四支箭上分别放置非零标量 \(\alpha,\beta,\gamma,\delta\)。这些无关系表示的同构类由
\[
\lambda=\frac{\gamma\alpha}{\delta\beta}\in k^\times
\]
唯一决定。因此 \(k\) 为无限域时，这个有限无环箭图已有无穷多个两两不同构的表示，各支箭的秩却都为一。若加入关系 \(ca-db=0\)，则上述全非零表示只剩 \(\lambda=1\) 的一个同构类。
\end{proposition}
\begin{proof}
在顶点 \(i\) 换基 \(g_i\in k^\times\)，两条从顶点 \(1\) 到顶点 \(4\) 的复合同时变为
\[
g_4(\gamma\alpha)g_1^{-1},\qquad
g_4(\delta\beta)g_1^{-1}.
\]
两者的比值因而不变。反过来，取
\[
g_1=1,\qquad g_2=\alpha^{-1},\qquad
g_3=\beta^{-1},\qquad g_4=(\gamma\alpha)^{-1},
\]
便把 \(a,b,c\) 上的标量都变为一，把 \(d\) 上的标量变为 \(\lambda^{-1}\)。所以同一比值的表示同构，不同比值的表示不同构。关系的矩阵方程是 \(\gamma\alpha=\delta\beta\)，恰好要求 \(\lambda=1\)。
\end{proof}
这里没有有向闭路径，但沿两条不同路径到达同一顶点时，两个复合之间的比例仍不能被顶点换基消去。下一节所讨论的有向直线没有这种分叉后重新汇合的情形，其表示可由各段复合的秩分类。
